\section{Recursive derivatives and derivative-component level sets}
\subsection{Representatives label estimates; they do not create information}
The affine gradient is constant throughout its source simplex. To turn the
cell gradients into another sampled field, assign a point to each retained
cell. The default is its centroid,
$\vect z_\sigma=(d+1)^{-1}\sum_j\vect x_j$, which stays inside the cell.
The code stores this point and the indices of its source vertices.
That provenance is needed to reason about contamination.

A circumcenter is an alternative. Subtracting squared distance equations gives
\begin{equation}
 2\matr E_\sigma(\vect c_\sigma-\vect x_0)
 =\begin{bmatrix}\norm{\vect e_1}^2&\cdots&\norm{\vect e_d}^2\end{bmatrix}^{\mathsf T}.
\end{equation}
For the isotropic quadratic $f=\norm{\vect x}^2/2$, the secant gradient equals
the exact gradient at the circumcenter. This explains why that representative
can look attractive. It is not generally true for an anisotropic Hessian.
An obtuse cell can place its circumcenter outside its convex hull: for
$(0,0),(2,0),(0.4,0.3)$ it is approximately $(1,-0.917)$, while the centroid is
$(0.8,0.1)$. Recursing on the former can extend support into an unmeasured
region. This optional mode is retained for experiments; duplicate
representatives cause an explicit error.

\subsection{Each recursion differentiates a newly interpolated field}
Let $V^{(0)}=f$. For order $k$, triangulate the support of $V^{(k-1)}$ and
apply \cref{eq:edgesolve} to every value component with the same edge matrix.
Place the resulting derivatives at new representatives. Ordered derivative
indices are retained; the number of components is $d^k$.

At order two, component-major flattening gives
\begin{equation}
 V^{(2)}=[\partial_1 g_1,\partial_2 g_1,\ldots,
          \partial_1 g_2,\partial_2 g_2,\ldots].
\end{equation}
Reshaping produces $H^{(2)}_{ij}=\partial_j g_i$ on the second-order support.
\glsadd{sym:recoveredhessian}
Its trace is the recovered Laplacian in the declared metric.
This matrix is the Jacobian of the reinterpolated gradient cloud. It is not
the classical Hessian of the original piecewise-affine scalar interpolant:
that Hessian is zero inside each original cell, with interface contributions
in a weak interpretation.

The raw recovered matrix need not be symmetric. Expose
\begin{equation}
 \matr H_{\mathrm{sym}}=\tfrac12(\matr H^{(2)}+
          (\matr H^{(2)})^{\mathsf T}),\qquad
 a_H=\norm{\matr H^{(2)}-(\matr H^{(2)})^{\mathsf T}}_{\mathrm F}.
 \label{eq:symmetry}
\end{equation}
Symmetrization is the closest symmetric matrix in Frobenius norm.
It removes an incompatible component without establishing the accuracy of
the remaining curvature. A small trace cannot exclude a saddle:
$\diag(1,-1)$ has zero trace and opposite curvatures. Eigenvalues of an
accurate symmetric Hessian can distinguish those directions; the trace alone
cannot.

An affine field is a useful limiting test: its recovered first gradient is
constant, and the next derivative is zero up to roundoff. Exact quadratic
Hessians require more. Secant gradients of a quadratic are not generally its
exact gradients at centroids; retriangulating those approximate values does
not repair that discrepancy. We test both an analytically supplied linear
vector gradient and the complete scalar recursion.

\subsection{What an edge intersection represents}
On the derivative mesh, consider component $v_c$ along an edge from
$\vect z_a$ to $\vect z_b$. For level $\ell$,
\begin{equation}
 t=\frac{\ell-v_{a,c}}{v_{b,c}-v_{a,c}},\qquad
 \vect p=(1-t)\vect z_a+t\vect z_b,\qquad 0\leq t\leq1.
 \label{eq:intersection}
\end{equation}
Every value component is interpolated with this same $t$. The selected
interpolated component should equal the requested level to numerical
precision, giving an independent test of the geometry.

The VICE backend prepares unique undirected edges and accepts lengths strictly
below $r$ times the RMS length of all unique edges; the default is $r=3$.
This filters long connections on the derivative mesh, not bad source cells
in the gradient solve. Equal-value edges are skipped: their intersection is
either empty or the whole edge, not one unique point. Shared-vertex cuts
remain once per incident edge. The result is an edge-intersection cloud,
not a connected contour mesh or a complete higher-dimensional isosurface.

By default, $M$ levels span $[\min v_c,\max v_c]$ uniformly. A quantile option
provides an explicit sensitivity experiment. Triangulation uses spatial
coordinates only. Field amplitude does not enter its distance metric, but
one very large derivative value can stretch the uniform level range and
leave few levels in the ordinary part of the field.

\begin{figure}[htbp]
\centering\includegraphics[width=\textwidth]{figures/generated/iso_geometry.pdf}
\caption{Derivative-component cuts on the same Himmelblau sample geometry,
before and after value perturbations. Each panel uses 15 uniformly spaced
levels over its own component range. Perturbed slopes change both that range
and cut locations. The dots are edge cuts, not connected contour lines.}
\label{fig:isogeometry}
\end{figure}

\subsection{Iso coverage is a counting-measure descriptor}
Pool the intersections over levels, retaining per-edge multiplicity, into
$\mathcal P_c$. At derivative support point $\vect z_i$, define
\begin{equation}
 s_{i,c}=\frac1{|\mathcal P_c|}\sum_{\vect p\in\mathcal P_c}
 \exp\!\left(-\frac{\norm{\vect p-\vect z_i}_2^2}{2h^2}\right),\qquad
 S_i=\left(\prod_{c=1}^{d^k}s_{i,c}\right)^{1/d^k}.
 \label{eq:coverage}
\end{equation}
\glsadd{sym:isocoverage}\glsadd{sym:kernelbandwidth}
The default $h$ is the RMS unique-edge length before filtering. The code
evaluates the mean over all intersections in memory-bounded chunks; it does
not silently substitute nearest neighbors or subsample the mesh.
Defined component scores lie in $[0,1]$, with Gaussian underflow possibly
producing zero. If any component has no usable cuts, the fused score is
undefined. Nearly constant components use the declared test
\begin{equation}
 \max v_c-\min v_c\leq10^{-10}\max(1,\max|v_c|).
\end{equation}
Reconsider this tolerance for very small physical derivative units.

The normalization by $|\mathcal P_c|$ makes this a kernel integral of an
empirical probability measure. It is not contour length or surface area:
edges and levels supply unequal geometric sampling, and vertex multiplicity
remains. For $d>2$, edge cuts additionally omit the interior of level-set
facets. Dense cuts indicate nearby derivative-level crossings, which can
come from genuine curvature, noise or contamination as well as a smoothly
resolved field.

If a smooth component varies approximately linearly and neighboring levels
differ by $\Delta\ell$, their hyperplanes have separation approximately
$\Delta\ell/\norm{\nabla v_c}$. Strong variation packs them closer.
Cut proximity can therefore encode variation while accuracy depends on noise
and geometry. As $h\to\infty$, every defined component score tends to one
and discrimination disappears. As $h\to0$, scores concentrate at sampled cuts.

\paragraph{The affine counterexample.}
For an exactly affine scalar field, all derivative components are constant.
The gradient is exact, yet its components yield no unique cuts. Assigning a
tiny score to missing cuts would classify an ideal case as unreliable.
The active code returns an undefined score and a
\texttt{constant\_component} status. Absence of derivative variation is not
absence of derivative information.

\subsection{Historical threshold and conservative search proposals}
For comparison with MATLAB, retain
\begin{equation}
 \tau=\mu_S+\lambda\sigma_S,\qquad
 \widehat y_i=\mathbb I[S_i<\tau].
 \label{eq:threshold}
\end{equation}
The standard deviation uses the sample convention. Increasing $\lambda$
raises $\tau$ and can only increase the number of low-score flags.
For skewed scores the threshold may exceed almost all values.
It is not a calibrated false-positive rule.

A nonzero first gradient supplies direction
$\vect d_i=\pm\vect g_i/\norm{\vect g_i}$ in the supplied coordinate metric.
The default step is one quarter of the distance from the representative to
its second-nearest original sample:
\begin{equation}
 \vect x_{i,\mathrm{trial}}=\vect z_i+\alpha_i\vect d_i.
 \label{eq:proposal}
\end{equation}
Zero gradients give zero steps. The sign selects ascent or descent.
The predicted linear gain magnitude is $\alpha_i\norm{\vect g_i}$, not a
guarantee of improvement. Proposals outside the original convex hull are
flagged. The hull fills holes and nonconvex forbidden regions; being inside
does not imply physical feasibility.

At $K>1$, the final coverage has a different support. Nearest-neighbor
transfer to first-order representatives returns both the descriptor and
transfer distance, even if the two support arrays have the same row count.
The default step is not scaled by uncalibrated coverage. A future acceptance
rule should evaluate a new measurement, or use independently justified error
bounds and a line search. A machine application must still impose current,
voltage, thermal and measurement-validity constraints.
